\section{Conclusion}
\label{sec:conclusion}

This paper presents a stochastic harmonic hosting capacity method that accounts for the uncertain angle of the harmonic currents injected by power-electronic interfaced units, and thereby for harmonic recombination. The angle uncertainty of a 6-pulse line-commutated converter is characterized through a harmonic state-space model and Monte Carlo simulation, and represented by one Beta-distributed germ per unit, which captures the strong dependence of its angles across the harmonic orders. A germ shared by all units of a converter type represents a common firing delay at little extra cost. The angles are referenced to the fundamental voltage at each unit's terminal, and their trigonometric functions are projected onto the polynomial chaos basis exactly by quadrature. The general problem is a non-convex quadratically constrained program; under the absolute equality fairness criterion it reduces to a second-order cone program with one pair of variables per harmonic, independent of the network size and the cardinality of the basis, which reduces exactly to the deterministic model when the angle ranges collapse.

The numerical illustration on the IEEE 30-bus system shows that harmonic recombination is limited when all units are of one converter type: the contributions arrive nearly in phase at the binding buses, the equivalent summation exponent stays close to one at the 5th and 7th harmonic, and the indicative exponents of IEC~61000-3-6 overestimate the hosting capacity by a factor of about two. Capacitive elements change the hosting capacity mainly through the network impedance and resonances, and add recombination only at the higher harmonics. Substantial recombination requires units with different mean angles, such as a mix of converter technologies: a balanced mix raises the 95\,\% hosting capacity by about one third and brings the equivalent exponent of the low-order harmonics close to~1.4, while the higher orders stay below the indicative~2.0. The deterministic angle is not a safe reference, and the value of the chance constraint lies in its gain over a robust bound.

Future work will replace the working assumptions on the phase-controlled converter by a harmonic state-space characterization of its firing delay, fit the trigonometric coefficients directly from such data, and extend the uncertainty to background distortion and network parameters.
